\section{La Mosca Soldado Negra (\textit{Hermetia illucens}) como solución biotecnológica}
La Hermetia illucens, o Mosca Soldado Negra (BSF, por sus siglas en inglés), se ha consolidado como una alternativa biotecnológica prometedora para la gestión de residuos orgánicos debido a su alta eficiencia en la biodegradación y conversión de materia orgánica en biomasa rica en proteínas y lípidos \parencite{barragan2017nutritional}. Este proceso no solo reduce el volumen de los desechos, sino que también minimiza las emisiones de gases de efecto invernadero (GEI) en comparación con métodos convencionales como el compostaje o la incineración. Además, sus subproductos, como el frass (larvicompost) y la biomasa larvaria, tienen aplicaciones en la producción de alimentos para animales, fertilizantes orgánicos y bioproductos industriales, fomentando así un modelo de economía circular \parencite{Diener2011,salamEffectDifferentEnvironmental2022,brunoValorizationOrganicWaste2025,koyunogluBiofuelProductionUtilizing2024,suryatiLauricAcidBlack2023,siddiquiEnhancingBioconversionRate2024,hamamAdvancesInsectIndustry2024,sarangiFoodWasteFood2024}.

Para optimizar la cría y el rendimiento de la BSF, es fundamental controlar con precisión variables ambientales como temperatura, humedad, pH y composición del sustrato, ya que estas afectan directamente la eficiencia del proceso de bioconversión y la calidad del producto final \parencite{salamEffectDifferentEnvironmental2022,belperioAssessingSubstrateUtilization2024,nayakHermetiaIllucensDiptera2024,goldEffectsRearingSystem2022}. Sin embargo, la ausencia de sistemas automatizados y herramientas avanzadas para el monitoreo en tiempo real limita la escalabilidad y eficiencia de esta tecnología en la gestión de residuos \parencite{surendraRethinkingOrganicWastes2020}.

Ante esta problemática, la integración de tecnologías emergentes, como el Internet de las Cosas (IoT) y algoritmos de aprendizaje automático (AAA), se presenta como una solución innovadora para optimizar el control de variables críticas, mejorar la trazabilidad del proceso y maximizar la eficiencia operativa. Esta convergencia tecnológica permitiría una gestión de residuos más eficiente y sostenible, alineada con los principios de la economía circular y la reducción de emisiones de GEI.